% Simulazione d'esame: prova del 01/02/1970, esercizio 1 -- riflessione di un urto obliquo su lastra a M = 1,3.
% Stessi numeri e figure dello svolgimento di riferimento:
% svolgimenti/1970-02-01_E01_riflessione_urto (calcoli.py e figure/).
% Compilare con: bash ../compila.sh 1970-02-01_E01_riflessione_urto  (A4 + orologio)

% Preambolo comune delle simulazioni d'esame (vedi GUIDA.md).
% Lo stesso sorgente si compila in due formati:
%   - A4 (svolgimento.pdf), da leggere su carta, PC o telefono;
%   - orologio (orologio.pdf): colonna stretta, letta dall'app dell'orologio
%     (garmin-latex), che taglia le pagine in strisce larghe quanto lo schermo.
% Il formato orologio si sceglie definendo \formatoorologio prima di \input
% (lo fa compila.sh). In quel formato il margine destro è volutamente largo:
% ciò che sporge oltre la colonna (equazioni lunghe, tabelle, figure) viene
% riconosciuto dall'app e mostrato rimpicciolito, con lo zoom.

\documentclass[11pt]{article}

\usepackage[T1]{fontenc}
\usepackage[utf8]{inputenc}
\usepackage[italian]{babel}
\usepackage{amsmath,amssymb}
\usepackage{array}
\usepackage{siunitx}
\usepackage{booktabs}
\usepackage{graphicx}
\usepackage{float}
\usepackage{caption}
\usepackage{enumitem}
\usepackage{ragged2e}
\usepackage{titlesec}
\usepackage{hyperref}

\hypersetup{hidelinks}
\setlength{\parindent}{0pt}
\setlength{\parskip}{4pt}
\allowdisplaybreaks
\sisetup{output-decimal-marker={,}, group-separator={\,}, per-mode=symbol,
         exponent-product=\cdot, range-phrase={\text{--}}}

% Notazione del formulario (come template/svolgimento.tex).
\newcommand{\Cp}{C^{+}}
\newcommand{\Cm}{C^{-}}
\newcommand{\Mr}{M_{r}}
\newcommand{\Ms}{M_{s}}

% Separatore fra relazioni nella stessa formula: \begin{gather*} A \ee B \end{gather*}.
% In A4 le relazioni stanno sulla stessa riga, nel formato orologio una sotto l'altra.
\ifdefined\formatoorologio \newcommand{\ee}{,\\} \else \newcommand{\ee}{,\qquad} \fi

% Titolo dell'esercizio: \esercizio{Prova del GG/MM/AAAA -- Esercizio n}{argomento}.
% Va anche nei metadati del PDF: l'app dell'orologio lo usa come titolo nell'elenco.
\newcommand{\esercizio}[2]{%
  \title{#1\\[2pt]\large #2}\author{}\date{}%
  \hypersetup{pdftitle={#1: #2}}}

\ifdefined\formatoorologio
  % Colonna di 4,6 cm: a 2,4 px/pt è larga 262 px sullo schermo da 390.
  \usepackage[paperwidth=12cm,paperheight=150cm,left=0.2cm,top=0.2cm,
              textwidth=4.6cm,textheight=149cm,nohead,nofoot]{geometry}
  \usepackage[scaled=1.0]{helvet}
  \usepackage{sfmath}
  \renewcommand{\familydefault}{\sfdefault}
  \pagestyle{empty}
  \makeatletter\let\ps@plain\ps@empty\makeatother   % anche la pagina del titolo
  % colonna dichiarata per l'app dell'orologio (deve coincidere con left e textwidth)
  \hypersetup{pdfkeywords={colonna=0.2cm+4.6cm}}
  \RaggedRight
  \setlength{\parskip}{5pt}
  \hfuzz=20cm          % ciò che sporge va nel margine destro, senza avvisi
  \hbadness=10000
  \titleformat*{\section}{\normalsize\bfseries\RaggedRight\hyphenpenalty=10000}
  \titleformat*{\subsection}{\normalsize\bfseries\RaggedRight}
  \setlist{leftmargin=*,itemsep=2pt}
  \captionsetup{font=small,justification=raggedright,singlelinecheck=false}
  % figure: larghe il doppio della colonna, così sporgono e l'app le mostra con lo zoom
  \newcommand{\figura}[3][0.7]{%
    \begin{figure}[H]\includegraphics[width=9cm]{#2}\caption{#3}\end{figure}}
\else
  \usepackage[a4paper,margin=2.5cm]{geometry}
  \newcommand{\figura}[3][0.7]{%
    \begin{figure}[H]\centering\includegraphics[width=#1\textwidth]{#2}\caption{#3}\end{figure}}
\fi

\graphicspath{{../../svolgimenti/1970-02-01_E01_riflessione_urto/figure/}}

\esercizio{Prova del 01/02/1970 -- Esercizio 1}{Riflessione di un urto obliquo su una lastra piana a $M=1{,}3$}

\begin{document}
\maketitle

\emph{Traccia.} Una corrente uniforme supersonica, stazionaria e piana a $M=1{,}3$ scorre lungo una lastra piana
su cui incide un urto obliquo. Definire la polare pressione--deviazione, determinare le diverse riflessioni possibili
(gas ideale) e descriverle nel piano fisico $(x,y)$. (Stessa traccia di B69-E01, 25/01/2023.)

\section*{Che cosa chiede e con quale modello}
La corrente 1 è parallela alla lastra ($\theta_1=0$, $p_1=1$). Un urto obliquo generato altrove, per esempio da un
cuneo sopra la lastra, la devia verso la lastra di un angolo $\theta_I$; la lastra impone di nuovo $\theta=0$, e il
modo in cui ci riesce è la riflessione. L'intensità dell'urto incidente non è data: la leggo come richiesta di
classificare le riflessioni possibili al variare di $\theta_I$ (o dell'angolo d'urto $\sigma_I$), a $M_1=1{,}3$
fissato. Assumo aria, $\gamma=1{,}4$.

Uso il modello piano, stazionario, non viscoso, con gas perfetto: trascuro lo strato limite sulla lastra e la sua
interazione con l'urto. Gli urti sono rettilinei e fra di essi gli stati sono uniformi, quindi ogni urto è un punto
della polare $p$--$\theta$ dello stato a monte. Il modello è iperbolico solo dove la corrente è supersonica: finché
le regioni 2 e 3 lo sono, la riflessione si decide localmente nel punto d'impatto; una regione subsonica invece
riceve informazioni da valle, e lì le sole polari non bastano. Il campo è esterno e nessuno impone una
contropressione, quindi prendo il ramo debole per l'incidente e, salvo discuterlo, anche per il riflesso. Numero gli
stati come d'uso: 1 a monte, 2 dopo l'incidente, 3 dopo il riflesso, 4 dietro il gambetto o dietro il tratto normale
dell'urto.

\section*{Che cosa mi aspetto prima dei numeri}
L'urto incidente viene dall'alto: è un'onda $\Cm$ di compressione, che devia la corrente di $-\theta_I$ verso la
lastra e porta lo stato 2 sul ramo debole della polare di 1. Dietro di esso la corrente punta contro la lastra, che
impone $\theta=0$: dalla lastra deve nascere un'onda $\Cp$ che ruoti la corrente di $+\theta_I$. Su una parete
un'onda si riflette con la stessa natura, quindi il riflesso è una compressione, cioè un urto, e la pressione sulla
lastra sale due volte: $p_3>p_2>p_1$.

Il Mach basso è il punto decisivo. Il riflesso deve dare la stessa deviazione $\theta_I$ partendo da $M_2<M_1$, cioè
con una polare più piccola: esiste solo se $\theta_I\le\theta_{\max}(M_2)$. Crescendo $\theta_I$, $M_2$ cala e con lui
$\theta_{\max}(M_2)$: ci sarà un angolo di distacco $\theta_D$ oltre il quale la riflessione regolare (RR) è
impossibile. Oltre $\theta_D$ dovrò verificare se esiste una riflessione di Mach (MR) a tre urti: incidente, riflesso
e gambetto quasi normale, con una linea di scorrimento dal punto triplo.

Le configurazioni candidate nel piano fisico sono tre:
\begin{itemize}
\item RR: due urti rettilinei che si incontrano sulla lastra; regioni uniformi 1, 2, 3 con $\theta_1=\theta_3=0$ e
  $\theta_2=-\theta_I$;
\item MR: punto triplo staccato dalla lastra, gambetto quasi normale fino alla lastra con regione 4 subsonica dietro,
  linea di scorrimento fra 3 e 4 con $p_3=p_4$, $\theta_3=\theta_4$;
\item nessuna riflessione: un solo urto, rettilineo lontano dalla lastra, che si incurva e la raggiunge
  perpendicolarmente, con una sacca subsonica 4 chiusa da una linea sonica.
\end{itemize}

Conto anche equazioni e incognite. Nella RR, noto 2, il riflesso ha un'incognita (il suo angolo) e una condizione,
$\theta_3=0$: si interseca la polare $\Cp$ di 2 con l'asse $\theta=0$. Una sola onda nasce dalla parete e una sola
condizione vi è imposta, come deve essere. Nella MR, noto 2, le incognite sono gli angoli del riflesso e del gambetto
(cioè $p_3,\theta_3$ e $p_4,\theta_4$ sulle rispettive polari) e le condizioni sono $p_3=p_4$, $\theta_3=\theta_4$:
si interseca la polare $\Cp$ di 2 con il ramo forte della polare di 1. Il bilancio è chiuso, ma la soluzione può
non esistere.

Tre scorciatoie mi faranno risparmiare calcoli:
\begin{itemize}
\item per sapere se la RR esiste basta confrontare $\theta_I$ con $\theta_{\max}(M_2)$, quindi $\theta_D$ è la radice
  di $\theta_{\max}(M_2(\theta_I))=\theta_I$;
\item il criterio di von Neumann (polari che si incontrano sull'asse nel punto d'urto normale di 1) richiede
  $p_3=p_{n1}$; poiché $p_3$ cresce con $\theta_I$, basta confrontare il suo massimo, a $\theta_D$, con $p_{n1}$;
\item per $\theta_I\to0$ le onde sono acustiche, ognuna con $\Delta p/p\simeq\gamma M^2\theta/\sqrt{M^2-1}$, e la
  riflessione su parete raddoppia il salto: $p_3-p_1\simeq2(p_2-p_1)$. Mi servirà come controllo.
\end{itemize}

\section*{La polare di $M_1=1{,}3$}
Attraverso un urto di angolo $\sigma$ la componente tangenziale della velocità si conserva e la normale ha il salto
dell'urto normale con $M_{n1}=M_1\sin\sigma$. Da $V_t$ costante e dalla continuità ho
$\tan(\sigma-\theta)/\tan\sigma=\rho_1/\rho_2=[(\gamma-1)M_{n1}^2+2]/[(\gamma+1)M_{n1}^2]$, da cui
\begin{gather*}
\tan\theta=2\cot\sigma\,\frac{M_1^2\sin^2\sigma-1}{M_1^2(\gamma+\cos2\sigma)+2}\ee
\frac{p_2}{p_1}=1+\frac{2\gamma}{\gamma+1}\left(M_1^2\sin^2\sigma-1\right)\ee
M_2=\frac{M_{n2}}{\sin(\sigma-\theta)}
\end{gather*}
con $M_{n2}^2=(1+\delta M_{n1}^2)/(\gamma M_{n1}^2-\delta)$. Facendo variare $\sigma$ da $\mu_1$ a \ang{90} la coppia
$(\theta,p)$ descrive la polare, con $\pm\theta$ per le due famiglie. I punti notevoli sono:
\begin{itemize}
\item onda nulla: $\sigma=\mu_1=\arcsin(1/M_1)=\ang{50.3}$, $p/p_1=1$;
\item massima deviazione ($d\theta/d\sigma=0$): $\theta_{\max}=\ang{6.66}$ a $\sigma=\ang{69.4}$, $p/p_1=1{,}56$;
\item punto sonico ($M_2=1$): $\theta_s=\ang{6.32}$ a $\sigma=\ang{65.1}$, $p/p_1=1{,}46$;
\item urto normale: $\theta=0$, $p_{n1}/p_1=1{,}805$, $M=0{,}786$.
\end{itemize}
Il ramo debole va da $\mu_1$ al punto sonico; fra punto sonico e $\theta_{\max}$ è ancora debole, ma con valle
subsonica. Tutta la fenomenologia sta quindi in meno di \ang{7}.

\section*{Riflessione regolare: un esempio a $\theta_I=\ang{2}$}
Per lo stato 2 uso la relazione $\theta$--$\sigma$--$M$ con $M_1$ e $\theta_I$, ramo debole:
\begin{gather*}
\sigma_I=\ang{53.5}\ee M_{n1}=1{,}045\ee p_2/p_1=1{,}107\ee M_2=1{,}22\ee \theta_2=-\ang{2}.
\end{gather*}
Controllo subito: $M_{n1}>1$ e $\Delta s>0$, ramo debole con $M_2>1$ (lo è per ogni $\theta_I<\theta_s(M_1)=\ang{6.32}$).
La polare di 2 arriva a $\theta_{\max}(M_2)=\ang{4.59}>\theta_I$: la RR esiste.

Per lo stato 3 uso la stessa relazione con $M_2$ e deviazione $+\theta_I$ (onda $\Cp$, da $-\theta_I$ a 0), ramo
debole. L'angolo $\sigma_R$ è misurato dalla direzione della corrente 2, quindi l'urto riflesso forma con la lastra
l'angolo $\varphi_R=\sigma_R-\theta_I$:
\begin{gather*}
\sigma_R=\ang{58.9}\ee \varphi_R=\ang{56.9}\ee M_{n2}=1{,}048\ee p_3/p_2=1{,}115\ee p_3/p_1=1{,}23\ee M_3=1{,}14.
\end{gather*}
Controllo: $M_{n2}>1$, $\Delta s>0$, $\theta_3=0$ soddisfatta, e il riflesso è una compressione come l'incidente, come
deve essere su una parete. Il riflesso è più intenso dell'incidente ($p_3/p_2=1{,}115>p_2/p_1=1{,}107$, perché
$M_2<M_1$) e non è speculare: $\varphi_R=\ang{56.9}\ne\sigma_I=\ang{53.5}$, la legge ottica non vale. Con la
scorciatoia acustica a $\theta_I=\ang{0.5}$: $p_2-p_1=0{,}0253$ (la teoria lineare dà $0{,}0249$) e
$p_3-p_1=0{,}0514\simeq2(p_2-p_1)=0{,}0505$; il calcolo è coerente.

La polare di 2 incontra l'asse $\theta=0$ anche sul ramo forte: $p_{3'}/p_1=1{,}73$, $M_{3'}=0{,}84$,
$\varphi_{R'}=\ang{82.6}$. Richiederebbe una contropressione imposta a valle, assente in un campo esterno, e la
scarto.

\figura[1.0]{p_theta.pdf}{Piano $p$--$\theta$: polare di 1 (blu), polare di 2 (rosso; a tratto continuo il ramo
$\Cp$ utile al riflesso). (a) RR: intersezioni 3 (debole) e 3$'$ (forte) con l'asse $\theta=0$. (b) Distacco: polare
di 2 tangente all'asse. (c) Oltre il distacco: la polare di 2 non raggiunge l'asse ed è interna alla polare di 1.
\label{fig:ptheta}}

Nel pannello (a) della figura~\ref{fig:ptheta} si legge la procedura: le polari partono da $(0,1)$, non
dall'origine; da 1 lungo il ramo debole della polare di 1 fino a 2, poi lungo la polare di 2 fino all'asse. Nel piano
fisico (figura~\ref{fig:campo}a) le linee di corrente piegano di \ang{2} verso la lastra attraversando l'incidente e
tornano orizzontali attraverso il riflesso, più inclinato sulla lastra dell'incidente.

\section*{Fin dove esiste la RR: distacco e criterio sonico}
Ripeto il calcolo al variare di $\theta_I$:

\begin{table}[H]
\centering\small
\begin{tabular}{cccccccc}
\toprule
$\theta_I$ & $\sigma_I$ & $p_2/p_1$ & $M_2$ & $\theta_{\max}(M_2)$ & $\varphi_R$ & $p_3/p_1$ & $M_3$ \\
\midrule
\ang{1}    & \ang{51.8} & 1,051 & 1,263 & \ang{5.64} & \ang{53.0} & 1,107 & 1,224 \\
\ang{2}    & \ang{53.5} & 1,107 & 1,224 & \ang{4.59} & \ang{56.9} & 1,233 & 1,140 \\
\ang{3}    & \ang{55.3} & 1,167 & 1,184 & \ang{3.52} & \ang{63.7} & 1,415 & 1,028 \\
$\theta_S=\ang{3.15}$ & \ang{55.6} & 1,176 & 1,178 & \ang{3.36} & \ang{65.8} & 1,462 & 1,000 \\
$\theta_D=\ang{3.25}$ & \ang{55.8} & 1,183 & 1,173 & \ang{3.25} & \ang{69.7} & 1,537 & 0,955 \\
\bottomrule
\end{tabular}
\caption{RR (riflesso debole) al variare della deviazione dell'urto incidente.}
\end{table}

$M_2$ diminuisce e la polare riflessa si restringe finché è tangente all'asse $\theta=0$: dalla radice di
$\theta_{\max}(M_2(\theta_I))=\theta_I$ ottengo
\begin{gather*}
\theta_D=\ang{3.25}\ee \sigma_I=\ang{55.8}\ee M_2=1{,}17\ee p_2/p_1=1{,}18,
\end{gather*}
con, al distacco, $p_3/p_1=1{,}54$ e $M_3=0{,}96$ (pannello (b) della figura~\ref{fig:ptheta}). Già un po' prima la
regione 3 diventa subsonica: $M_3=1$ per $\theta_S=\ang{3.15}$ ($p_3/p_1=1{,}46$). Qui c'è il controllo critico:
fra $\theta_S$ e $\theta_D$ la regione 3 è subsonica ($M_3\ge0{,}96$) e il modello iperbolico vale solo vicino al
punto di riflessione.

\section*{Criterio di von Neumann}
Il massimo di $p_3$ sulla RR è al distacco: $p_3/p_1=1{,}54<p_{n1}/p_1=1{,}805$. La pressione dell'urto normale di 1
non è mai raggiunta, quindi le due polari non si incontrano mai sull'asse $\theta=0$ e l'angolo di von Neumann non
esiste a $M_1=1{,}3$. Manca anche il dominio di doppia soluzione: sotto $\theta_D$ c'è solo la RR, senza isteresi.

\figura[0.62]{p_parete.pdf}{Pressioni dietro l'incidente e a parete per la RR, al variare di $\theta_I$; in
punteggiato la pressione dell'urto normale di 1.\label{fig:pparete}}

La figura~\ref{fig:pparete} lo mostra: il ramo debole e quello forte della RR si congiungono a $\theta_D$, dove la RR
finisce, e $p_3$ resta sempre sotto $1{,}805$. Le pressioni stanno fra 1 e 1,8, quindi basta la scala lineare.

\section*{Oltre il distacco: la MR non esiste}
Per $\theta_D<\theta_I<\theta_s(M_1)$ cerco l'intersezione della polare $\Cp$ di 2 con il ramo forte della polare di 1.
Prendo $\theta_I=\ang{5}$:
\begin{gather*}
\sigma_I=\ang{60.0}\ee M_2=1{,}09\ee p_2/p_1=1{,}31\ee \theta_{\max}(M_2)=\ang{1.31}.
\end{gather*}
La polare di 2 occupa solo $-\ang{6.31}\le\theta\le-\ang{3.69}$ e arriva al massimo a $p/p_1=1{,}60$ (riflesso
normale), mentre il ramo forte della polare di 1 a $\theta=-\ang{5}$ è a $p/p_1=1{,}73$. La polare riflessa è tutta
interna alla polare dell'incidente (pannello (c) della figura~\ref{fig:ptheta}) e la tocca solo in 2; lo stesso
accade per ogni $\theta_I$ dell'intervallo. La teoria dei tre urti non ha soluzione: a $M_1=1{,}3$ non c'è MR.
Controllo la coerenza: $\theta_D=\ang{3.25}<\theta_s(M_1)=\ang{6.32}<\theta_{\max}(M_1)=\ang{6.66}$, e l'assenza di
$\theta_N$ è coerente con l'assenza di soluzioni a tre urti.

Per $\theta_s(M_1)<\theta_I<\theta_{\max}(M_1)$ la regione 2 è già subsonica e non può esistere un urto riflesso; per
$\theta_I>\theta_{\max}(M_1)$ l'urto incidente stesso non può essere attaccato al corpo che lo genera.

Che cosa succede allora oltre $\theta_D$? La polare riflessa non interseca né l'asse delle pressioni né la polare
dell'incidente: non c'è riflessione. L'urto incidente si incurva avvicinandosi alla lastra e la raggiunge
perpendicolarmente, così che la condizione $\theta=0$ a parete è soddisfatta dal tratto normale. Lungo l'urto curvo
lo stato a valle percorre la polare di 1, da 2 (tratto rettilineo, lontano dalla lastra) fino all'urto normale
($p_4/p_1=1{,}805$, $M_4=0{,}79$) al piede. Dietro il tratto quasi normale c'è una sacca subsonica, e l'urto curvo
dà entropia diversa a ogni linea di corrente: il campo a valle è rotazionale.

\figura[1.0]{campo.pdf}{Piano fisico. (a) RR a $\theta_I=\ang{2}$, angoli e linee di corrente in scala. (b) Oltre il
distacco: urto incidente che si incurva e incide perpendicolarmente la lastra, con sacca subsonica 4
(qualitativo).\label{fig:campo}}

Nella figura~\ref{fig:campo}(b) l'urto si raccorda alla lastra con un tratto quasi normale, senza urto riflesso né
punto triplo; a valle la corrente subsonica 4 riaccelera e si raccorda con la regione 2 attraverso la linea sonica.

\section*{Risposta: le riflessioni possibili a $M_1=1{,}3$}
\begin{table}[H]
\centering\small
\begin{tabular}{p{0.3\linewidth}p{0.42\linewidth}l}
\toprule
$\theta_I$ & Configurazione & Note \\
\midrule
$0<\theta_I<\theta_S=\ang{3.15}$ & RR, regioni 2 e 3 supersoniche & $p_3/p_1\le1{,}46$ \\
$\theta_S<\theta_I\le\theta_D=\ang{3.25}$ & RR con regione 3 subsonica & $p_3/p_1\le1{,}54$ \\
$\theta_D<\theta_I<\theta_s(M_1)=\ang{6.32}$ & nessuna riflessione, urto curvo & niente MR \\
$\theta_s(M_1)<\theta_I<\theta_{\max}(M_1)=\ang{6.66}$ & regione 2 subsonica, idem & \\
\bottomrule
\end{tabular}
\caption{Riflessioni possibili ($\gamma=1{,}4$); $\theta_N$ non esiste.}
\end{table}

Esempio di RR a $\theta_I=\ang{2}$: $\sigma_I=\ang{53.5}$, $p_2/p_1=1{,}107$, $M_2=1{,}22$; $\varphi_R=\ang{56.9}$,
$p_3/p_1=1{,}23$, $M_3=1{,}14$.

\figura[0.62]{diagramma_M_theta.pdf}{Piano $\theta_I$--$M_1$: deviazione massima e sonica di 1, criteri di distacco
$\theta_D$ e di von Neumann $\theta_N$; tratteggiata la retta $M_1=1{,}3$.\label{fig:diagr}}

La figura~\ref{fig:diagr} colloca il caso nel quadro generale: $\theta_N$ esiste solo per $M_1\gtrsim2{,}2$, dove si
apre la regione di doppia soluzione fra $\theta_N$ e $\theta_D$; a $M_1=1{,}3$ la retta verticale incontra solo
$\theta_D=\ang{3.25}$ e poi, a pochi gradi, $\theta_s$ e $\theta_{\max}$.

\section*{Discussione}
\begin{itemize}
\item A Mach basso tutto sta in pochi gradi: basta un incidente debole ($p_2/p_1\approx1{,}18$) per rendere impossibile
  la RR, perché il riflesso parte da $M_2\approx1{,}17$, la cui polare ammette meno di \ang{3.3} di deviazione.
\item La RR esiste fino a $\theta_D$, ma da $\theta_S$ la regione 3 è subsonica: le perturbazioni di valle (per esempio
  la spalla del cuneo che genera l'incidente) possono risalire fino al punto di riflessione, quindi la transizione
  reale avviene fra $\theta_S$ e $\theta_D$, qui appena \ang{0.1}.
\item La MR a tre urti richiede che la polare riflessa esca dalla polare incidente: a $M_1=1{,}3$ ne resta sempre
  dentro, quindi oltre il distacco non c'è riflessione e l'urto si incurva fino a incidere perpendicolarmente la
  lastra. Le misure reali in questo regime mostrano un urto che si raccorda con continuità alla parete, con una
  debole compressione riflessa che si dissolve, invece di un punto triplo netto.
\item Non c'è isteresi: senza $\theta_N$ la configurazione dipende solo da $\theta_I$ e non dal verso in cui lo si
  fa variare, a differenza di quanto accade per $M_1\gtrsim2{,}2$.
\item Il riflesso non segue la legge ottica ($\varphi_R\ne\sigma_I$): è più inclinato e più intenso perché deve dare
  la stessa deviazione partendo da un Mach minore.
\item Lo strato limite sulla lastra, che ho trascurato, a questi Mach interagisce fortemente con l'urto (la sua parte
  subsonica trasmette la sovrapressione a monte) e può separare: è il limite principale del modello non viscoso.
\end{itemize}

\end{document}
